\documentclass[11pt]{article}
\usepackage[a4paper,margin=24mm]{geometry}
\usepackage{fontspec}
\setmainfont{TeX Gyre Pagella}
\usepackage[UTF8]{ctex}
\usepackage{amsmath,amssymb,amsthm,mathtools}
\usepackage{xcolor}
\usepackage[colorlinks=true,linkcolor=blue!50!black,citecolor=blue!50!black,urlcolor=blue!50!black]{hyperref}
\hypersetup{pdftitle={非边界局部化商的局部幂零部分（修正版）}}
\newcommand{\C}{\mathbb C}
\newcommand{\Z}{\mathbb Z}
\newcommand{\g}{\widehat{\mathfrak{sl}}_2}
\newcommand{\slc}{\mathfrak{sl}_2}
\newcommand{\ad}{\operatorname{ad}}
\newcommand{\Span}{\operatorname{span}_{\C}}
\newtheorem{theorem}{定理}
\newtheorem{lemma}[theorem]{引理}
\newtheorem{proposition}[theorem]{命题}
\newtheorem{corollary}[theorem]{推论}
\setlength{\emergencystretch}{2em}
\title{非边界局部化商的局部幂零部分（修正版）\\
  \large 修正后的 E-作用与原初递推；撤回权滤过结论}
\author{进展报告（2026-09-22 审计后修正）}
\date{2026 年 9 月 22 日}
\begin{document}
\maketitle

\begin{abstract}
本文是 9 月 22 日 $G_c$ 结构笔记的修正版。独立审计
（\texttt{audits/2026-09-22/}）给出了显式反例，否定：旧的 E-作用公式、原初
空间描述 $P_\nu\cong\{X:T^{\nu+1}(X)u=0\}$、$\mu=-2$ 时 $P_0=\C w_1$ 的
断言、以及权尾部 $G^{\ge\nu}$ 是仿射子模的断言。本版记录：经全分量 $A_r$
修正的 E-作用、原初向量的修正三角递推、$P_0$ 中的无限族 $Z^nw_1$、含降权
映射 $\ell_j:P_\nu\to P_{\nu-2}$ 的完整投影组，并撤回权滤过的全部推论。
不受影响的既有结论：边界同构、深区单性判据、整数参数的唯一极大子模与
simple head、代数中心约化的相容性。$G_c$ 的内部分类（单性、socle、长度）
为开放问题。
\end{abstract}

\section{定义、记号与假设}

\paragraph{仿射代数与诱导模。}
沿用统一笔记：$\g=(\mathfrak{sl}_2\otimes\C[t,t^{-1}])\oplus\C K$，括号
\[
 [e_i,f_j]=h_{i+j}+i\delta_{i+j,0}K,\quad [h_i,h_j]=2i\delta_{i+j,0}K,\quad
 [h_i,e_j]=2e_{i+j},\quad [h_i,f_j]=-2f_{i+j},
\]
$K$ 中心，$[e_i,e_j]=[f_i,f_j]=0$。固定 $a,b\in\Z$、$a+b\ge0$、
$L=a+b+1\ge1$，$S_{a,b}=\Span\{K,h_i\ (i\ge0),e_i\ (i>a),f_j\ (j>b)\}$。
特征 $\chi$ 由 $\kappa=\chi(K)$ 与 $\lambda_i=\chi(h_i)$（$0\le i\le L$，
$i>L$ 时 $\lambda_i=0$）决定；$M_{a,b}(\chi)=U\otimes_{U(S_{a,b})}\C_\chi$，
$u=1\otimes1_\chi$。中心水平记为 $\kappa$（不再用 $c$），$c$ 专指局部化
指标。

\paragraph{深区局部化。}
取 $c<b$ 且 $c\le-a-1$（即 $d_*:=b-c\ge L$），$F=f_c$。谱流
$\sigma_{-c}$ 将 $\mathcal T_FM_{a,b}(\chi)$ 等同于
$Q=\mathcal T_{f_0}M_{A,B}(\chi')$，其中 $A=a+c\le-1$，$B=d_*\ge L$，
$\mu:=\chi'(h_0)=\lambda_0-c\kappa$。记 $E=e_0$，$H=h_0$，
$w_m=f_0^{-m}u+M$；则 $[E,f_0]=H$，$Eu=0$，$Hu=\mu u$。

\paragraph{PBW 记号。}
$\mathcal X$ 为字
$X=f_{j_1}\cdots f_{j_p}e_{i_1}\cdots e_{i_q}h_{r_1}\cdots h_{r_s}$
（$j_t\le B,\ j_t\ne0,\ i_t\le A,\ r_u<0$）的集合（含空字）。
$\{F^{-m}Xu:m\ge1,\ X\in\mathcal X\}$ 是 $Q$ 的基；$\deg X=2q-2p$，
$H(F^{-m}Xu)=(\mu+2m+\deg X)F^{-m}Xu$。令 $V_0$ 为 $\{Xu:X\in\mathcal X\}$
的张成（$F$-自由部分），则 $M=\bigoplus_{r\ge0}f_0^rV_0$。

\paragraph{局部幂零部分。}
$G:=\{v\in Q:E^nv=0\ \text{对某个}\ n\ge1\}$ 是 $\g$-子模；$G$ 的每个
非零 $\g$-子模都含 $P_\nu:=\ker E\cap G\cap Q_\nu$ 中的非零向量。浅区
$-a\le c<b$ 中 $E=e_{-c}$ 在 $Q$ 上单射，故 $G=0$；只有深区带有非零局部
幂零部分。

\section{修正后的 E-作用}

\begin{lemma}\label{lem:Eaction}
对 $x\in V_0$ 定义线性算子 $A_r:V_0\to V_0$：
\[
 Ex=\sum_{r\ge0}f_0^r\,A_rx,
\]
对每个 $x$ 为有限和，$A_0=\pi_0\circ\ad(e_0)$ 即旧的 $T$。若
$Hx=(\mu+\delta)x$ 且 $m\ge1$，则在 $Q$ 中
\begin{equation}\label{eq:Eaction}
 E\bigl(f_0^{-m}x\bigr)
 =\sum_{r=0}^{m-1}f_0^{-(m-r)}A_rx
 -m(\mu+\delta+m+1)f_0^{-m-1}x.
\end{equation}
\end{lemma}

\begin{proof}
由三重恒等式 $Ef_0^{-m}=f_0^{-m}E-mf_0^{-m-1}H-m(m+1)f_0^{-m-2}f_0$ 与
$Ex=\sum_rf_0^rA_rx$；$f_0^{-m}f_0^rA_rx$ 在 $Q$ 中当 $0\le r<m$ 时存活为
$f_0^{-(m-r)}A_rx$，$r\ge m$ 时落回 $M$。
\end{proof}

\begin{corollary}\label{cor:counterexample}
仅保留 $A_0$ 的旧公式错误。取 $X=f_{-1}f_1$，则
$EXu=(f_1h_{-1}+\lambda_1f_{-1}-2f_0)u$，故
\[
 E\bigl(f_0^{-2}f_{-1}f_1u\bigr)
 =f_0^{-2}(f_1h_{-1}+\lambda_1f_{-1})u
 -2w_1-2(\mu-1)f_0^{-3}f_{-1}f_1u,
\]
旧公式漏掉非零项 $-2w_1$（\texttt{work/check\_audit\_regressions.py}
机器核验）。
\end{corollary}

\section{原初空间}

\begin{lemma}\label{lem:pole}
若 $0\ne v\in Q_\nu$ 且 $Ev=0$，则 $\nu\in\Z_{\ge0}$，且 $v$ 的最高 pole
恰为 $\nu+1$。
\end{lemma}

\begin{proof}
写 $v=\sum_{m=1}^{N}f_0^{-m}x_m$，$x_m\in V_0$。在 pole $N+1$ 只有
\eqref{eq:Eaction} 的对角项存活，迫使 $N=\nu+1$。
\end{proof}

\begin{theorem}\label{thm:P}
非零向量 $v=\sum_{m=1}^{\nu+1}f_0^{-m}x_m\in Q_\nu$ 为原初，当且仅当其各层
满足三角系统
\begin{equation}\label{eq:recursion}
 \sum_{m=k}^{\nu+1}A_{m-k}x_m=(k-1)(\nu-k+2)x_{k-1},
 \qquad 1\le k\le\nu+1,\quad x_0=0,
\end{equation}
且 $k=1$ 处的相容条件为 $\sum_{m=1}^{\nu+1}A_{m-1}x_m=0$。$P_\nu$ 是该线性
系统的解空间，一般不是单个 monomial 的集合；旧描述
$P_\nu\cong\{X:T^{\nu+1}(X)u=0\}$ 错误。
\end{theorem}

\begin{proof}
用 \eqref{eq:Eaction} 比较 $Ev=0$ 中 $f_0^{-k}$ 的系数。$k=\nu+1$ 时右端
为 $\nu x_\nu=A_0x_{\nu+1}$，故顶层向下逐层确定；$k=1$ 即剩余相容条件。
$\nu=1$ 时系统为 $x_1=A_0x_2$ 与 $(A_0^2+A_1)x_2=0$。
\end{proof}

\begin{corollary}\label{cor:Zcounter}
旧判据确有反例。元素
\[
 Z=h_{-1}^2+4f_{-1}e_{-1}+2h_{-2}
 =h_{-1}^2+2(e_{-1}f_{-1}+f_{-1}e_{-1})
\]
与 $E$、$f_0$、$H$ 对易（不必在整个 $\g$ 中中心）。$\mu=-1$ 时
$p=\lambda_1f_0^{-1}u+f_0^{-2}f_1u\in P_1$，且 $Zp\in P_1$ 的各层为
$x_1=(\lambda_1Z-4h_{-1})u$、$x_2=(f_1Z-4(\kappa+1)f_{-1})u$，满足
$A_0^2x_2=8e_{-1}u\ne0$；缺失项 $A_1x_2=-8e_{-1}u$ 恢复相容性（机器核验）。
\end{corollary}

\begin{corollary}\label{cor:P0}
$\mu=-2$ 时 $\dim P_0=\infty$：向量 $Z^nw_1=f_0^{-1}Z^nu+M$（$n\ge0$）
原初且线性无关，其最高 PBW 次数为 $2n$，主符号
$(h_{-1}^2+4f_{-1}e_{-1})^n\ne0$。撤回旧断言 $P_0=\C w_1$。
\end{corollary}

\section{\texorpdfstring{$\slc$}{sl2}-分解与投影}

\begin{theorem}\label{thm:sl2}
$G$ 的每个向量落在有限维 $\slc(E,H,f_0)$-子模中，且作为横向 $\slc$-模
\[
 G\;\cong\;\bigoplus_{\nu\in\Z_{\ge0}}P_\nu\otimes L(\nu),
\]
其中 $L(\nu)$ 为 $(\nu+1)$ 维不可约模。此分解正确并保留；它不是
$\g$-子模分解。
\end{theorem}

\begin{proof}
$E$ 与 $f_0$ 在 $G$ 上局部幂零，故每个向量生成有限维 $\slc$-模；完全可约
是 Weyl 定理，原初权为非负整数。
\end{proof}

全部投影的原理：current 模态在横向 $\slc$ 的伴随作用下构成 $L(2)$，而
$L(2)\otimes L(\nu)\cong L(\nu+2)\oplus L(\nu)\oplus L(\nu-2)$
（$\nu\ge2$）。

\begin{proposition}\label{prop:proj}
设 $p\in P_\nu$ 在 pole $\nu+1$ 的顶层为 $X$。
\begin{enumerate}
\item[(i)] $i\le A$：$e_ip\in P_{\nu+2}$；$\nu=0$ 时顶层为 $-2f_iu$
（例：$\mu=-2$ 时 $e_{-1}w_1=f_0^{-1}e_{-1}u-f_0^{-2}h_{-1}u
-2f_0^{-3}f_{-1}u$）；一般情形最高 pole $\nu+3$ 由 $[e_i,f_0^{-m}]$ 的修正层
承载。
\item[(ii)] $r<0$、$\nu\ge1$：$h_rp+\tfrac{2}{\nu+2}f_0e_rp\in P_\nu$，顶层
为 $\tfrac{\nu}{\nu+2}(-h_rX+\tfrac1\nu f_rA_0X)$；$\nu=0$ 时
$h_rp=-f_0e_rp$。
\item[(iii)] $j\ge1$：$E^2f_jp=-2e_jp$，$f_jp$ 有 $P_{\nu+2}$ 与 $P_\nu$
两个分量。
\item[(iv)] $\nu\ge2$、任意 $j$：降权投影
\begin{equation}\label{eq:lowering}
 \ell_j(p)=f_jp-\tfrac1\nu f_0h_jp
 -\tfrac{1}{\nu(\nu+1)}f_0^2e_jp\;\in\;P_{\nu-2},
\end{equation}
即上述张量积的 $L(\nu-2)$-分量。它确实非零：$\mu=-2$ 时对
$p_2=e_{-1}w_1\in P_2$ 有
\[
 \ell_1(p_2)=f_0^{-1}\Bigl(\tfrac13f_1e_{-1}+\tfrac{\lambda_1}{6}h_{-1}
 +\tfrac{2\kappa-2}{3}\Bigr)u+M\in P_0\setminus\{0\}.
\]
\end{enumerate}
\end{proposition}

\begin{proof}
(ii)(iii) 同前用 $[E,h_r]=-2e_r$、$[f_0,e_r]=-h_r$、$[E,f_j]=h_j$；(iv) 对
\eqref{eq:lowering} 右端施加 $E$，并用 $p$ 的块上
$Ef_0^s=s(\nu-s+1)f_0^{s-1}$ 验证。非零例机器核验。
\end{proof}

\section{子模结构：撤回与开放状态}

\paragraph{权尾部不是仿射子模。}
旧稿声称 $G^{\ge\nu}:=\bigoplus_{\nu'\ge\nu}P_{\nu'}\otimes L(\nu')$ 是
$\g$-子模。这是错误的。
\begin{enumerate}
\item[(i)] 中心障碍：$\mu=-2$、$\kappa\ne0$ 时取 $p=w_1\in L(0)$。
$h_{\pm1}p=-f_0e_{\pm1}p$ 落在顶层权为 $2$ 的块中，故 $h_1p$、$h_{-1}p$
都在 $G^{\ge2}$ 内；若 $G^{\ge2}$ 是子模，则
$[h_1,h_{-1}]p=2\kappa p\in G^{\ge2}$，与 $p\in L(0)$、$\kappa\ne0$
矛盾（机器核验）。
\item[(ii)] 降权通道：$\mu=-2$ 时 $\ell_1(e_{-1}w_1)\in P_0\setminus\{0\}$
说明模态把 $P_2$ 映进 $P_0$，故 $\kappa=0$ 时尾部同样不稳定。
\end{enumerate}
因此在水平 $\kappa\ne0$ 时 $Q$ 没有非零平凡仿射商模：$K$ 作用为 $\kappa$，
而所有 $h$-模态作用为零会迫使 $[h_1,h_{-1}]=2K$ 作用为零。

\paragraph{撤回的断言。}
以下由无效滤过推出的结论撤回为已证结果：$G$ 永不单；
$\operatorname{soc}(G)=0$ 与 $\operatorname{soc}(Q)=0$；$G$ 合成长度无限；
商 $G^{\ge\nu}/G^{\ge\nu+1}$ 单并构成 head。其否命题同样未证：$G_c$ 的
内部分类为开放问题。

\paragraph{保留的结论。}
边界同构；深区单性判据 $\mathcal T_{f_c}M_{a,b}$ 单 $\iff
\lambda_0-c\kappa\notin\Z$（在 $\lambda_L\ne0$、$\kappa\ne-2$ 下）；整数参数
的唯一极大子模 $G_c$、simple head $Q_c/G_c$ 与不分裂扩张；代数中心约化的
相容性；$G_c$ 的横向 $\slc$-分解；本文的修正递推与命题 \ref{prop:proj} 的
投影。正确的 $G_c$ 子模分析必须同时追踪 $L(\nu+2)$、$L(\nu)$、$L(\nu-2)$
三个通道；旧滤过弃用。

\section{验证状态}

审计 \texttt{audits/2026-09-22/}（报告、154 项独立精确检查、结果、日志）已
纳入仓库。四类反例由 \texttt{work/check\_audit\_regressions.py} 复现
（24 项）：E-作用残差 $-2w_1$；截断 $T$ 残差 $-4e_{-2}u$；$P_0$ 的 $Z$-族
与 $\mu=-1$ 的修正递推例；障碍 $[h_1,h_{-1}]w_1=2\kappa w_1$；非零
$\ell_1(p_2)$。这些是佐证而非形式证明；子模分类的开放状态应在任何项目摘要
中如实反映。

\begin{thebibliography}{9}
\bibitem{FGXZ}
V.~Futorny, X.~Guo, Y.~Xue and K.~Zhao,
\emph{Smooth representations of affine Kac--Moody algebras},
arXiv:2404.03855v2 (2024), Theorem~1.1.
\url{https://arxiv.org/abs/2404.03855v2}.

\bibitem{Audit}
2026-09-22 独立审计：\texttt{audits/2026-09-22/audit\_report\_cn.md} 与
\texttt{audits/2026-09-22/independent\_pbw\_audit.py}（154 项精确检查）。

\bibitem{Unified}
项目笔记：\emph{Boundary and Nonboundary Localization Quotients of Smooth
Affine Modules}（统一笔记，2026 年 9 月）及 2026-09-15 附录。
\end{thebibliography}
\end{document}
